% !TeX program = pdflatex
% Компиляция: bash build_paper.sh  (pdflatex прогоняется до схождения
% перекрёстных ссылок; двух проходов не хватает, нужно три)
\documentclass[11pt,a4paper]{article}

% cmap ДОЛЖЕН идти до fontenc: он встраивает в PDF таблицы ToUnicode, без
% которых кириллица не извлекается из готового файла — текст нельзя ни найти
% поиском, ни скопировать, ни разобрать системам индексации (РИНЦ, Crossref).
\usepackage{cmap}
\usepackage[T2A]{fontenc}
\usepackage[utf8]{inputenc}
\usepackage[english,russian]{babel}
\usepackage{amsmath,amssymb,amsthm}
\usepackage{booktabs}
\usepackage{graphicx}
% Рисунки генерирует analysis.py и кладёт в figs/.
\graphicspath{{figs/}{./}}
\usepackage[margin=25mm]{geometry}
% xurl ДО hyperref: разрешает переносить длинные URL в любом месте. Без него
% ссылки вроде t5k.org/mersenne/heuristic.html вылезают за правое поле.
\usepackage{xurl}
\usepackage[colorlinks=true,allcolors=blue]{hyperref}
% Запас на растяжение межсловных пробелов: в русском тексте с большим числом
% неразрывной формульной вёрстки без него остаются переполненные строки.
\emergencystretch=3em
\usepackage{xcolor}
\usepackage[expansion=false]{microtype}
% Подписи в принятом отечественном виде: «Рис. 1. Название», «Таблица 1. Название»
% (вместо «Рис. 1: Название»).
\usepackage{caption}
\captionsetup{labelsep=period}
% п. 4.1.5 ГОСТ Р 7.0.7-2021: подписи к иллюстративному материалу приводят на
% языке текста статьи и, как правило, повторяют на английском языке.
% Номер берётся из счётчика, уже увеличенного \caption, поэтому русская и
% английская подписи всегда несут один и тот же номер.
\newcommand{\encaptab}[1]{\\[3pt]\foreignlanguage{english}{%
  \textbf{Table~\thetable.} #1}}
\newcommand{\encapfig}[1]{\\[3pt]\foreignlanguage{english}{%
  \textbf{Figure~\thefigure.} #1}}

\theoremstyle{plain}
\newtheorem{proposition}{Предложение}
\newtheorem{corollary}{Следствие}
\theoremstyle{remark}
\newtheorem{remark}{Замечание}

\newcommand{\Rb}{R_b}
\newcommand{\eg}{e^{\gamma}}
\newcommand{\emg}{e^{-\gamma}}
\DeclareMathOperator{\Gam}{Gamma}

\title{Схема наблюдения в статистике обобщённых репьюнитных\\
простых: калибровка объединённой проверки константы\\
Ленстры --- Померанса --- Вагстаффа}
\author{Тимофей Докучаев\thanks{ORCID: \url{https://orcid.org/0009-0006-0510-5225}.
Контактный адрес: \texttt{timofei93@yandex.ru}}\\
\small Независимый исследователь, Иркутск, Россия}
% Дата фиксирована: 	oday менялась бы при каждой пересборке.
\date{27 августа 2026~г.}

\begin{document}

% Тип статьи (п. 4.4) и индекс УДК (п. 4.5) ГОСТ Р 7.0.7-2021:
% приводят в начале статьи отдельными строками слева.
\noindent Научная статья\\
\noindent УДК 511.3:519.234
\medskip

\maketitle

\begin{abstract}
Эвристика Ленстры --- Померанса --- Вагстаффа (ЛПВ) предсказывает, что
$\#\{n\le x:\Rb(n)\ 	ext{просто}\}\sim(\eg/\ln b)\ln x$. С 1993~года эту универсальность проверяют через
<<нормированную частоту встречаемости>>, получаемую регрессией $\ln n_k$ на~$k$.

Во-первых, подозрение о смещении этой частоты правым цензурированием
\emph{неверно}: и концевая форма, и наклон МНК несмещены точно
(предложение~\ref{prop:unbiased}); смещена лишь обратная величина --- неравенство
Йенсена. Дефект иной: варианты оценки расходятся до $29{,}5\%$ --- в $1{,}7$ раза
больше медианного отклонения от $\emg$, --- и ярлык <<удачливости>> основания
переворачивается от смены формы.

Во-вторых, точное условное правдоподобие вместо регрессии объединяет $B=15$
оснований в $M=218$ событий против $51$ у одних простых Мерсенна. Главное: пивот
$\kappa S\sim\Gam(M,1)$ точен лишь при наблюдении <<до $N$-го события>>, а данные
получены при фиксированном фронте со случайным числом находок. При реальной схеме номинальный $5\%$
критерий отвергает в $16{,}2\%$ случаев, а верной объединённой оценкой оказывается
$(M-B)/S$, а не $(M-1)/S$: условное правдоподобие стоит одного события \emph{на
каждое основание}. Отсюда $\hat\kappa=1{,}838$ при $95\%$ интервале
$[1{,}61;\,2{,}13]$ и $p=0{,}64$ против $\eg\approx1{,}781$; сообщавшееся
десятипроцентное превышение над ЛПВ есть артефакт схемы наблюдения.

В-третьих, выбор между двумя оценками определяется способом сбора данных, а не
самими данными, и стоит $7\%$; отсюда требование публиковать фронты $L_b$.
Наконец, критерий однородности ($\Lambda=7{,}40$, $p=0{,}92$) различает лишь
разброс $\kappa_b$ свыше $\pm34\%$, а объединённая выборка --- отклонения от
$\eg$ свыше $21\%$.

\medskip
\noindent\textbf{Ключевые слова:} репьюнит-простые, схема наблюдения, калибровка,
пуассоновский процесс, условное правдоподобие, гипотеза Ленстры --- Померанса --- Вагстаффа

\medskip
\noindent\textbf{MSC 2020:} Primary 11Y11; Secondary 11N05, 62F03, 62B05, 11N80.

\medskip
\noindent\textbf{Благодарности:} автор благодарит проект Great Internet Mersenne
Prime Search (GIMPS) за предоставление данных о простых числах Мерсенна через
открытый отчёт PrimeNet (\url{https://www.mersenne.org/}); значения индексов
Мерсенна получены оттуда 27~августа 2026~года. Автор признателен участникам
OEIS, которые много лет поддерживают последовательности обобщённых репьюнитных
простых.

\medskip
\noindent\textbf{Финансирование:} исследование выполнено без внешнего
финансирования.
\end{abstract}

\begin{otherlanguage}{english}

% Англоязычный блок: заглавие, автор, аффилиация, аннотация, ключевые слова.
% Требуется журналами перечня ВАК и всеми индексируемыми в Scopus/WoS.
\begin{center}
{\large\bfseries The observation scheme in the statistics of generalized\\
repunit primes: calibrating the pooled test of the\\
Lenstra--Pomerance--Wagstaff constant}

\medskip
{\bfseries Timofei Dokuchaev}

\smallskip
% п. 4.9 ГОСТ Р 7.0.7-2021: организация, город, страна, электронный адрес
% (без слова «e-mail») и ORCID разделяют запятыми; точку в конце не ставят.
{\small Independent researcher, Irkutsk, Russia,\\
\texttt{timofei93@yandex.ru}, \url{https://orcid.org/0009-0006-0510-5225}}
\end{center}

\begin{abstract}
\noindent\textbf{Abstract.} The Lenstra--Pomerance--Wagstaff (LPW) heuristic
predicts $\#\{n\le x: (b^n-1)/(b-1)\ \text{prime}\}\sim(e^{\gamma}/\ln b)\ln x$.
Since 1993 it has been tested via a ``normalised occurrence frequency'' from
regressing $\ln n_k$ on $k$.

First, the suspicion that right-censoring biases this frequency downwards is
\emph{false}: both the endpoint form and the OLS slope are exactly unbiased; only
the reciprocal is biased --- Jensen's inequality. The genuine defect is that
competing forms of the estimator disagree by up to $29.5\%$ on identical data, and the
``lucky base'' label flips between them.

Second, an exact conditional likelihood pools $B=15$ bases into $M=218$ events
against $51$ for Mersenne primes alone. The main result concerns the pivot:
it is exact only under inverse sampling, whereas the data arise from a
fixed frontier with a random number of finds. Under the latter the nominal $5\%$
test rejects in $16.2\%$ of cases, and the correct pooled estimator is $(M-B)/S$,
not $(M-1)/S$: the conditional likelihood costs one event \emph{per base}. This
yields $\hat\kappa=1.838$, $95\%$ interval $[1.61,\,2.13]$, $p=0.64$ against
$e^{\gamma}\approx1.781$; the reported ten-percent excess is an artefact of the
scheme.

Third, the choice between the two estimators is settled by how the data were
collected, not by the data, and costs $7\%$; hence our call to publish the
frontiers $L_b$. Finally, the homogeneity test ($\Lambda=7.40$, $p=0.92$) detects
only a spread of $\kappa_b$ above $\pm34\%$, the pooled sample only deviations
above $21\%$.

\medskip
\noindent\textbf{Keywords:} repunit primes, observation scheme, calibration, Poisson
process, conditional likelihood, Lenstra--Pomerance--Wagstaff conjecture

\medskip
\noindent\textbf{Acknowledgments:} the author thanks the Great Internet Mersenne
Prime Search (GIMPS) project for providing the Mersenne prime data through the
open PrimeNet report (\url{https://www.mersenne.org/}); the Mersenne exponents
were retrieved on 27 August 2026. The author is grateful to the OEIS contributors
who have maintained the generalized repunit prime sequences for many years.

\medskip
\noindent\textbf{Funding:} this research received no external funding.
\end{abstract}
\end{otherlanguage}

\section{Введение}

Для целого основания $b\ge2$ обобщённый репьюнит есть
\begin{equation}\label{eq:Rb}
\Rb(n)=\frac{b^{n}-1}{b-1},
\end{equation}
причём $\Rb(n)$ может быть простым только при простом~$n$. При $b=2$ это числа Мерсенна
$M_p=2^{p}-1$. Опираясь на эвристики Вагстаффа~\cite{wagstaff1983} и
Померанса~\cite{pomerance1981}, гипотеза ЛПВ утверждает, что
\begin{equation}\label{eq:LPW}
\pi_b(x)=\#\{n\le x:\ \Rb(n)\ \text{простое}\}\ \sim\ \frac{\eg}{\ln b}\,\ln x,
\end{equation}
где $\gamma$~--- постоянная Эйлера~--- Маскерони; подробный вывод эвристики для
$b=2$ изложен в справочнике PrimePages~\cite{t5k}. Основание входит
в~\eqref{eq:LPW} лишь через множитель $1/\ln b$: константа $\eg\approx1{,}7811$
объявляется универсальной.
Дабнер~\cite{dubner1993} предложил проверять эту универсальность эмпирически по разным
основаниям, а Бурделэ~\cite{bourdelais2009} с тех пор ведёт текущие оценки нормированной
частоты, публикуя их в комментариях к записям OEIS~\cite{oeis}. Родственные вопросы для
одной лишь последовательности Мерсенна недавно обсуждались в~\cite{wagstaff2025}.
Близка по духу и работа~\cite{dominguez2026}, где у показателей Мерсенна ищут
вторичную структуру, связанную с разложением $p-1$; методически она устроена так же,
как наш~\S\ref{sec:struct} --- стратифицированное условное правдоподобие плюс
проверка значимости, --- но проверяет структуру \emph{внутри} одного основания,
тогда как мы проверяем её \emph{между} основаниями.

Вокруг этих оценок сложился интерпретационный фольклор. Основание~13 описывается как
<<удачливое>>, основание~7 --- как <<самое неудачливое>>; отмечается, что нормированные
частоты лежат ниже предполагаемого предела $\emg\approx0{,}5615$ и медленно сходятся
к нему снизу. Настоящая работа разбирает эти утверждения, но приходит к выводам,
отличным от тех, которых естественно ожидать.

\subsection{Вклад работы}

\begin{enumerate}
\item\textbf{Употребительная оценка частоты не смещена.} Схема, при которой интервал
подгонки заканчивается наибольшим известным индексом, действительно обусловливает
оценку тем, что последнее наблюдение является событием, и достаточная статистика
распределена как $\Gam(N,\lambda_b)$ (предложение~\ref{prop:gamma}). Но отсюда
\emph{не} следует смещение нормированной частоты: и концевая форма
$D/(N\ln b)$, и наклон МНК несмещены точно (предложение~\ref{prop:unbiased},
таблица~\ref{tab:mc}). Смещена только обратная величина
$\hat\kappa=1/\hat G$, в $N/(N-1)$ раз, и это неравенство Йенсена, а не цензурирование.
Несмещённые оценки $\kappa$ и $1/\kappa$ не могут быть взаимно обратными; всякая
<<поправка>>, чинящая одну, портит другую (\S\ref{sec:notbias}).

\item\textbf{Настоящий дефект --- неустойчивость к выбору формы оценки.}
В обращении несколько численно различных вариантов. На одних и тех же данных они
расходятся до $29{,}5\%$ (основание~18) --- в $1{,}7$ раза больше медианного
отклонения оценок от $\emg$ ($17{,}6\%$), которое и толкуют содержательно; у двух
оснований из пятнадцати выбор формы весит больше, чем всё их отклонение от
предсказания ЛПВ. Основание~18 при концевой форме --- третье
по <<неудачливости>> из пятнадцати, а при МНК --- десятое, ниже медианы. Ярлык
описывает выбор оценки, а не свойства простых (\S\ref{sec:instability}).

\item\textbf{Объединение оснований и его цена.} Условное правдоподобие
(\S\ref{sec:condlik}) не требует знания недокументированных границ поиска $L_b$ и
позволяет объединить $B=15$ оснований в $M=218$ событий против $51$ у одних
Мерсенна. Точное нулевое распределение $\Gam(M,1)$ делает симуляции ненужными.

\item\textbf{Главный результат: пивот не откалиброван под реальную схему наблюдения.}
Условное правдоподобие точно при схеме <<наблюдаем до $N$-го события>>, где $N$ задано
планом. Данные же получены при фиксированном фронте: поиск покрыл $(t_0,L_b]$, а число
находок случайно. При этой схеме номинальный $5\%$ критерий отвергает
в $16{,}2\%$ случаев, а правильной объединённой оценкой оказывается $(M-B)/S$
(предложение~\ref{prop:schemeB}). Это даёт $\hat\kappa=1{,}838$ при $p=0{,}64$:
превышение над ЛПВ исчезает (\S\ref{sec:scheme}).

\item\textbf{Выбор схемы данными не решается.} Оценки $(M-1)/S$ и $(M-B)/S$ построены
на одних и тех же достаточных статистиках $(N_b,D_b)$ и различаются на $7\%$;
что из них верно, определяется не данными, а тем, как они собирались. Попытка обойти
выбор через заведомо несмещённые по-основанию куски проваливается
(\S\ref{sec:noescape}). Единственный радикальный выход --- публиковать фронты $L_b$.

\item\textbf{Удачливые основания: чего критерий не может.} Критерий отношения
правдоподобий даёт $\Lambda=7{,}40$ при $14$ степенях свободы, $p=0{,}92$. Но его
мощность такова, что он различает лишь разброс $\kappa_b$ порядка $\pm34\%$
(\S\ref{sec:homog}). Корректный вывод --- <<базозависимые частоты данными не
требуются>>, а не <<их не существует>>.

\item\textbf{Явный расчёт мощности.} Объединённая выборка различает отклонения от
$\eg$ не менее $21{,}2\%$; одна последовательность Мерсенна --- не менее $49{,}8\%$
(\S\ref{sec:power}).
\end{enumerate}

В~\S\ref{sec:sameerror} разобран более частный, но поучительный сюжет: наивная регрессия
счётной функции Мерсенна порождает на вид значимый отрицательный свободный член,
который оказывается артефактом обусловливания. Мы включаем его как разобранный пример
той же ошибки в ином обличье.

\section{Модель и данные}\label{sec:data}

\subsection{Пуассоновская модель}
Положим $t=\ln n$. Эвристика~\eqref{eq:LPW} равносильна утверждению, что индексы
обобщённых репьюнитных простых по основанию~$b$ образуют однородный пуассоновский
точечный процесс на оси~$t$ с постоянной интенсивностью
\begin{equation}\label{eq:lam}
\lambda_b=\frac{\kappa}{\ln b},\qquad \text{ЛПВ:}\ \ \kappa=\eg=1{,}781072\ldots
\end{equation}
Универсальность --- это утверждение, что единственный параметр $\kappa$ управляет всеми
основаниями. Пуассоновская идеализация здесь --- рабочее допущение, а не теорема:
даже для обычных простых пуассоновский закон в коротких промежутках доказан лишь
условно, при гипотезах Харди --- Литтлвуда, и известно, что он разрушается при росте
масштаба~\cite{jha2026}. Мы принимаем~\eqref{eq:lam} как модель и всюду рассматриваем
$\kappa$ как оцениваемую величину; это сводит вопрос <<верна ли ЛПВ?>> к однопараметрическому и
делает основания объединяемыми.

\subsection{Допустимые основания}\label{sec:bases}
Основания, являющиеся точными степенями, исключаются: при $b=c^{m}$, $m>1$, число
$\Rb(n)$ допускает алгебраическую факторизацию, и процесс не имеет предполагаемого вида.
Поэтому основания $4,8,9,16$ опущены. Внутри объявленного диапазона отбор
не зависит от результатов: берутся \emph{все} основания $2\le b\le20$, не являющиеся
точными степенями, что даёт $B=15$; ни одно основание не исключено и не добавлено
по величине его оценки. Каждое из них имеет не менее восьми известных индексов
(минимум достигается при $b=18$: восемь членов A133857).

Подчеркнём, однако, что сама верхняя граница $b\le20$ --- условная круглая величина,
а не следствие критерия <<не менее восьми индексов>>: этот критерий выполняется и
дальше (A127996 для $b=21$ содержит восемь членов, A127997 для $b=22$~--- девять) и
впервые нарушается лишь при $b=23$ (A204940, шесть членов). Чтобы граница не выглядела
подобранной, мы проверили её влияние прямо. Расширение набора до всех оснований
$b\le26$, не являющихся точными степенями (добавляются $b=21,22,23,24,26$; $b=25=5^2$
исключается алгебраически), даёт $B=20$, $M=253$ и
\[
  \frac{M-1}{S}=1{,}9776,\qquad \frac{M-B}{S}=1{,}8285,
\]
то есть сдвигает основную оценку~\eqref{eq:resB} на $0{,}5\%$ --- вчетверо меньше
цены выбора схемы наблюдения (\S\ref{sec:scheme}). Критерий однородности на расширенном
наборе даёт $\Lambda=10{,}27$ при $19$ степенях свободы, $p=0{,}95$. Ни один вывод
работы от границы $b\le20$ не зависит; мы сохраняем её как заранее объявленную.

\subsection{Данные и их статус}
В таблице~\ref{tab:data} перечислены использованные последовательности,
полученные 27~августа 2026~года (сверка --- приложение~\ref{app:data}).
Три оговорки существенны.

\paragraph{Вероятно простые числа.} Для больших индексов большинство элементов ---
вероятно простые, установленные тестом Ферма или сильным тестом Ферма после пробного
деления. Счёт смещается вверх лишь в случае, если псевдопростое Ферма было принято за
простое, что для чисел такого размера практически не встречается. Чувствительность
проверена в~\S\ref{sec:robust}.

\paragraph{Фронты проверки.} Для $b=2$ проект GIMPS~\cite{gimps} публикует два фронта:
все показатели ниже $X_{\mathrm{ver}}=81\,648\,221$ проверены дважды, что даёт ровно $50$
простых Мерсенна, тогда как все показатели ниже $X_{\mathrm{ft}}=141\,308\,443$ проверены
хотя бы однократно, что даёт $52$. Два наибольших показателя предварительны: между ними
остаются непроверенные кандидаты. Анализ~\S\ref{sec:sameerror}, где фронт используется
явно, ограничен $50$ проверенными показателями. Объединённый анализ к этому выбору
нечувствителен (\S\ref{sec:robust}).

Оба фронта --- движущиеся величины, поэтому фиксируем момент считывания: отчёт
PrimeNet Milestones, 27~августа 2026~г. Ни один вывод от точного значения
$X_{\mathrm{ver}}$ не зависит: достаточно того, что оно лежит между
$M_{50}=77\,232\,917$ и $M_{51}=82\,589\,933$, а это верно для всех значений отчёта
за 2026~год.

\paragraph{Статус записи A000043.} Раздел DATA этой записи содержит $50$ членов и
обрывается на $77\,232\,917$: порядок следующих показателей OEIS считает недоказанным
до завершения двойной проверки. Показатели $82\,589\,933$ и $136\,279\,841$ известны и
документированы GIMPS (и упомянуты в комментариях к A000043), но формально не входят в
DATA. Мы включаем их, поскольку статистика~\eqref{eq:pivot} использует лишь положение
последнего события, а не его порядковый номер; в~\S\ref{sec:robust} проверено, что
ограничение основания~2 пятьюдесятью дважды проверенными показателями вывода не меняет.

\paragraph{Границы поиска при $b>2$.} Для прочих оснований сопоставимого фронта не
публикуется. Именно эта лакуна --- предмет~\S\ref{sec:scheme} и, как мы полагаем,
главное препятствие к количественной проверке ЛПВ.

\begin{table}[t]
\centering
\caption{Использованные последовательности. $N_b$ --- число событий строго правее
$t_0=\ln n_{\min}$, то есть на единицу меньше числа известных простых индексов:
наименьший индекс задаёт начало отсчёта, а не событие. Источник --- записи OEIS,
получено 27.08.2026.
\encaptab{Sequences used. $N_b$ is the number of events strictly to the right of
$t_0=\ln n_{\min}$, that is, one less than the number of known prime indices: the
smallest index sets the origin, not an event. Source: OEIS entries, retrieved
27.08.2026.}}
\label{tab:data}
\small
\begin{tabular}{rlrrrrr}
\toprule
$b$ & OEIS & простых & $N_b$ & $n_{\min}$ & $n_{\max}$ & $D_b$\\
\midrule
2 & A000043 & 52 & 51 & 2 & 136\,279\,841 & 18{,}0371\\
3 & A028491 & 23 & 22 & 3 & 8\,530\,117 & 14{,}8605\\
5 & A004061 & 19 & 18 & 3 & 3\,300\,593 & 13{,}9110\\
6 & A004062 & 17 & 16 & 2 & 3\,360\,347 & 14{,}3344\\
7 & A004063 & 10 & 9 & 5 & 1\,264\,699 & 12{,}4409\\
10 & A004023 & 11 & 10 & 2 & 8\,177\,207 & 15{,}2237\\
11 & A005808 & 13 & 12 & 17 & 1\,868\,983 & 11{,}6077\\
12 & A004064 & 14 & 13 & 2 & 769\,543 & 12{,}8604\\
13 & A016054 & 13 & 12 & 5 & 1\,503\,503 & 12{,}6139\\
14 & A006032 & 11 & 10 & 3 & 1\,724\,417 & 13{,}2618\\
15 & A006033 & 10 & 9 & 3 & 639\,833 & 12{,}2704\\
17 & A006034 & 12 & 11 & 3 & 1\,990\,523 & 13{,}4053\\
18 & A133857 & 8 & 7 & 2 & 1\,270\,141 & 13{,}3615\\
19 & A006035 & 11 & 10 & 19 & 209\,359 & 9{,}3074\\
20 & A127995 & 9 & 8 & 3 & 984\,349 & 12{,}7011\\
\midrule
\multicolumn{3}{l}{всего} & $M=218$ & & & $S=110{,}4423$\\
\bottomrule
\end{tabular}
\end{table}

\section{Что смещено, а что нет}\label{sec:notbias}

\subsection{Употребительные оценки}
Дабнер~\cite{dubner1993} и Бурделэ~\cite{bourdelais2009} характеризуют основание~$b$
нормированной частотой встречаемости, получаемой подгонкой прямой к зависимости
$\ln n_k$ от~$k$ с последующим делением на $\ln b$. Согласно~\eqref{eq:lam} ожидаемый
наклон равен $1/\lambda_b=\ln b/\kappa$, так что нормированная частота оценивает
\begin{equation}\label{eq:G}
G_b=\frac{1}{\kappa},\qquad\text{ЛПВ:}\ \ G_b=\emg=0{,}5614595\ldots
\end{equation}
В обращении несколько вариантов --- обычный МНК со свободным членом, регрессия через
начало координат и концевая форма $\hat G_b=D_b/(N_b\ln b)$, где $D_b=t_{N_b}-t_0$.
Все они разделяют структурную особенность: интервал подгонки заканчивается на наибольшем
известном индексе.

\subsection{Точное распределение}

\begin{proposition}[Условный закон правого конца]\label{prop:gamma}
Пусть события образуют однородный пуассоновский процесс интенсивности $\lambda$ на
$(t_0,\infty)$, и пусть $t_1<\dots<t_N$ --- его первые $N$ событий после $t_0$.
Тогда $D:=t_N-t_0$ имеет распределение $\Gam(N,\lambda)$ с плотностью
$\lambda^{N}d^{N-1}e^{-\lambda d}/(N-1)!$. При условии $D=d$ предшествующие $N-1$ событий
распределены как порядковые статистики $N-1$ независимых равномерных величин на $(t_0,d)$.
\end{proposition}

\begin{proof}
Времена ожидания $t_1-t_0,\dots,t_N-t_{N-1}$ независимы и распределены по закону
$\mathrm{Exp}(\lambda)$, а сумма $N$ таких величин распределена как $\Gam(N,\lambda)$.
Условная равномерность есть свойство порядковых статистик пуассоновского процесса,
применённое к $(t_0,d)$.
\end{proof}

\begin{corollary}[Смещение обратной величины]\label{cor:kappa}
Оценка максимального правдоподобия для $\lambda$ по одному лишь $D$ равна
$\hat\lambda=N/D$, причём
\begin{equation}\label{eq:bias}
\mathbb E[\hat\lambda]=N\,\mathbb E[D^{-1}]=\frac{N}{N-1}\,\lambda,
\end{equation}
так как $\mathbb E[D^{-1}]=\lambda/(N-1)$ для $D\sim\Gam(N,\lambda)$. Оценка
$\tilde\lambda=(N-1)/D$ в точности несмещена.
\end{corollary}

\subsection{Чего из следствия \ref{cor:kappa} не следует}

Естественно заключить, что раз $\hat\lambda$ завышена в $N/(N-1)$ раз, то частота
$\hat G_b$ занижена во столько же. Это неверно.

\begin{proposition}[Несмещённость частоты]\label{prop:unbiased}
В условиях предложения~\ref{prop:gamma}:
\begin{enumerate}
\item[\rm(i)] концевая оценка несмещена точно:
$\mathbb E\!\left[\dfrac{D}{N\ln b}\right]=\dfrac{1}{\lambda\ln b}=G_b$;
\item[\rm(ii)] наклон МНК регрессии $t_k$ на $k$ несмещён точно:
$\mathbb E[\hat\beta]=1/\lambda$;
\item[\rm(iii)] следовательно, <<исправленная>> величина
$\tilde G_b=\frac{N}{N-1}\hat G_b$ смещена \emph{вверх} ровно в $N/(N-1)$ раз.
\end{enumerate}
\end{proposition}

\begin{proof}
(i) $\mathbb E[D]=N/\lambda$ для $D\sim\Gam(N,\lambda)$, откуда результат немедленно.
(ii) Наклон МНК есть линейная статистика $\hat\beta=\sum_k w_k t_k$ с весами
$w_k=(k-\bar k)/\sum_j(j-\bar k)^2$, для которых $\sum_k w_k=0$ и $\sum_k w_k k=1$.
Поскольку $\mathbb E[t_k]=t_0+k/\lambda$, получаем
$\mathbb E[\hat\beta]=\sum_k w_k\,(t_0+k/\lambda)=1/\lambda$.
(iii) непосредственно из (i).
\end{proof}

Содержательный смысл прост: несмещённые оценки $\kappa$ и $1/\kappa$ не могут быть
взаимно обратными, поскольку $x\mapsto1/x$ строго выпукло. Несмещённая пара --- это
$(\hat G_b,\ \tilde\kappa_b)$, где $\tilde\kappa_b=(N_b-1)\ln b/D_b$; но
$\tilde\kappa_b\ne1/\hat G_b$. Всякая единая <<поправка на цензурирование>>,
исправляющая одну величину, портит другую. Таблица~\ref{tab:mc} подтверждает
предложение~\ref{prop:unbiased} симуляцией.

\begin{table}[t]
\centering
\caption{Проверка предложений методом Монте-Карло, $200\,000$ реализаций на строку
(зерно генератора фиксировано, см.\ \texttt{analysis.py}).
Целевое значение всех столбцов, кроме второго, равно $1{,}0000$.
\encaptab{Monte Carlo check of the propositions, $200\,000$ replications per row
(generator seed fixed, see \texttt{analysis.py}). The target value of every column
except the second is $1.0000$.}}
\label{tab:mc}
\small
\begin{tabular}{rrrrrrr}
\toprule
$N$ & $\mathbb E[\hat\lambda]/\lambda$ & $N/(N-1)$ & $\mathbb E[\tilde\lambda]/\lambda$
& $\mathbb E[\hat G]\lambda$ & $\mathbb E[\tilde G]\lambda$ & $\mathbb E[\hat\beta]\lambda$\\
\midrule
7 & 1{,}1665 & 1{,}1667 & 0{,}9998 & 1{,}0002 & 1{,}1669 & 1{,}0002\\
10 & 1{,}1115 & 1{,}1111 & 1{,}0004 & 0{,}9998 & 1{,}1109 & 0{,}9997\\
16 & 1{,}0672 & 1{,}0667 & 1{,}0005 & 0{,}9994 & 1{,}0660 & 0{,}9994\\
22 & 1{,}0481 & 1{,}0476 & 1{,}0004 & 0{,}9996 & 1{,}0472 & 0{,}9997\\
51 & 1{,}0202 & 1{,}0200 & 1{,}0002 & 0{,}9999 & 1{,}0199 & 0{,}9998\\
\bottomrule
\end{tabular}
\end{table}

\begin{remark}
Отрицательный результат этого параграфа стоит подчеркнуть, поскольку он противоречит
интуиции: систематическое отклонение опубликованных частот вниз от $\emg$
\emph{не} объясняется схемой подгонки. Его источник --- либо реальное превышение
$\kappa$ над $\eg$, либо, как показывает \S\ref{sec:scheme}, схема наблюдения,
но не арифметика оценки.
\end{remark}

\subsection{Неустойчивость к выбору формы оценки}\label{sec:instability}

Несмещённость не означает согласованности. В таблице~\ref{tab:bases} приведены
концевая форма $\hat G_b$ и оценка по МНК $\hat G_b^{\mathrm{OLS}}$ на одних и тех же
данных. Расхождения достигают $29{,}5\%$ ($b=18$), $15{,}9\%$ ($b=13$) и $10{,}3\%$
($b=11$).

Сопоставим их с тем, что в литературе толкуют содержательно, --- с отклонениями
самих оценок от $\emg$. Последние по пятнадцати основаниям имеют медиану $17{,}6\%$
(от $5{,}6\%$ при $b=20$ до $43{,}7\%$ при $b=19$). Наибольшее расхождение форм,
$29{,}5\%$, превышает эту медиану в $1{,}7$ раза. В среднем выбор формы, разумеется,
дешевле: медиана расхождений равна $4{,}1\%$, то есть около четверти медианного
отклонения. Но у двух оснований из пятнадцати ($b=18$ и $b=13$) выбор формы весит
больше, чем всё их отклонение от предсказания ЛПВ, --- и именно эти основания
фольклор и обсуждает.

Следствие для фольклора прямое. Основание~18 при концевой форме имеет
$\hat G_{18}=0{,}6604$ --- третье по величине из пятнадцати (после $b=7$ с $0{,}7104$
и $b=10$ с $0{,}6612$), то есть уверенно входит в тройку <<самых неудачливых>>; при МНК
$\hat G_{18}^{\mathrm{OLS}}=0{,}4656$ --- десятое место из пятнадцати, ниже медианы
$0{,}4904$, то есть оно уже <<удачливое>>. Ярлык переворачивается без единого нового
простого числа, от одной лишь смены формы оценки: основание перемещается с третьего
места на десятое.

К этому добавляется вторая, техническая беда МНК: его стандартные ошибки в этой задаче
неверны. Остатки суть отклонения порядковых статистик и сильно автокоррелированы;
для $b=2$ мы получаем $\rho_1=0{,}825$, а ковариационная оценка Ньюи --- Уэста с восемью
лагами увеличивает стандартную ошибку свободного члена с $0{,}447$ до $0{,}732$
(\S\ref{sec:sameerror}). Опубликованные точечные значения не сопровождаются интервалами,
и это, а не смещение, --- главная методологическая претензия к текущей практике.

\subsection{Условное правдоподобие}\label{sec:condlik}

Предложение~\ref{prop:gamma} даёт больше, чем разбор смещения. Полагая $t_0=\ln n_{\min}$
и приняв $(N_b,D_b)$ за данные по основанию~$b$, получаем логарифмическое правдоподобие
\begin{equation}\label{eq:loglik}
\ell(\kappa)=\sum_b\Bigl[N_b\ln\frac{\kappa}{\ln b}-\frac{\kappa}{\ln b}D_b\Bigr]+\mathrm{const},
\end{equation}
распределение которого не зависит от границы поиска $L_b$. Это существенно практически:
$L_b$ хорошо документирована лишь при $b=2$. Дифференцируя~\eqref{eq:loglik} и полагая
$S=\sum_b D_b/\ln b$, получаем
\begin{equation}\label{eq:mle}
\hat\kappa=\frac{M}{S},\qquad \tilde\kappa=\frac{M-1}{S},\qquad M=\sum_b N_b .
\end{equation}
Так как $\kappa S=\sum_b\lambda_bD_b$ есть сумма независимых величин $\Gam(N_b,1)$, то
\begin{equation}\label{eq:pivot}
\kappa S\sim\Gam(M,1),
\end{equation}
то есть точная центральная статистика; доверительные интервалы и $p$-значения получаются
в замкнутой форме. \emph{Ровно это утверждение и требует ревизии} --- см.~\S\ref{sec:scheme}.

\begin{remark}[Независимость по основаниям]
Объединение предполагает независимость процессов по разным основаниям. Строго она не
выполняется на малых индексах: семнадцать значений индекса встречаются более чем в одном
основании, причём $n=3$ --- в девяти, поскольку $\Rb(3)=b^2+b+1$ просто для многих~$b$.
Все такие совпадения лежат при $n\le317$. Существеннее, впрочем, не совпадение индексов
(числа $\Rb(n)$ при разных $b$ различны), а возможная общность алгебраических делителей;
и то и другое затрагивает лишь малые~$n$, тогда как $D_b$ определяется верхним концом
диапазона. Раздел~\ref{sec:robust} подтверждает, что отбрасывание всех индексов ниже~$10^3$
вывод не меняет.
\end{remark}

\section{Две схемы наблюдения}\label{sec:scheme}

\subsection{В чём различие}

Пивот~\eqref{eq:pivot} выведен в предположении, что $N_b$ задано планом наблюдения:
мы смотрим на процесс, пока не наберём $N_b$ событий, и останавливаемся. Назовём это
\emph{схемой~A} (обратная выборка, цензурирование типа~II).

Данные получены иначе. Для каждого основания вычислительный поиск покрыл интервал
$(t_0,\ln L_b]$ целиком, и найденными оказались все события, попавшие внутрь; число
находок $N_b$ --- случайная величина, а фиксирован фронт $L_b$. Назовём это
\emph{схемой~B}.

Различие не косметическое. При схеме~B, условно на $N_b=n$, положения событий суть
$n$ независимых равномерных величин на $(t_0,\ln L_b]$, поэтому условное распределение
$D_b$ зависит от $L_b$ и \emph{не} зависит от $\lambda_b$: при фиксированном $N_b$
величина $D_b$ информативна о фронте, а не об интенсивности. Правдоподобие~\eqref{eq:loglik}
при схеме~B не является правдоподобием.

Заметим также, что \S\ref{sec:sameerror} настоящей работы, где для $b=2$ фронт известен
и используется явно, построен именно на схеме~B. Обе схемы в одной работе применять
нельзя; ниже мы принимаем схему~B как соответствующую способу получения данных
и калибруем под неё результаты схемы~A.

\subsection{Правильная объединённая оценка при схеме B}

Смещение имеет замкнутую форму, и она поучительна.

\begin{proposition}\label{prop:schemeB}
Пусть при схеме~B поиск по основанию~$b$ покрыл $(t_0,\ln L_b]$, ширина окна есть
$W_b=\ln L_b-t_0$, а $\delta_b=\ln L_b-t_{N_b}$ --- расстояние от фронта до последнего
найденного события (при $N_b=0$ полагаем $\delta_b=W_b$). Тогда
\begin{equation}\label{eq:delta}
\Pr(\delta_b>x)=e^{-\lambda_b x}\ \ (0\le x<W_b),\qquad
\Pr(\delta_b=W_b)=e^{-\lambda_b W_b},
\end{equation}
то есть $\delta_b$ есть экспоненциальная величина с параметром $\lambda_b$,
усечённая на $W_b$, и
\begin{equation}\label{eq:Edelta}
\mathbb E[\delta_b]=\frac{1-e^{-\lambda_b W_b}}{\lambda_b}.
\end{equation}
Полная экспозиция есть $S_{\mathrm{full}}=S+\sum_b\delta_b/\ln b=\sum_b W_b/\ln b$,
причём $\mathbb E[M]=\kappa S_{\mathrm{full}}$ тождественно, а
\begin{equation}\label{eq:EBd}
\mathbb E\Bigl[\kappa\sum_b\delta_b/\ln b\Bigr]=B-\sum_b e^{-\lambda_b W_b}.
\end{equation}
Приравнивая $\kappa S_{\mathrm{full}}$ к $M$ и отбрасывая экспоненциально малый
остаток, получаем моментную оценку
\begin{equation}\label{eq:MBS}
\hat\kappa_{\mathrm B}=\frac{M-B}{S}.
\end{equation}
\end{proposition}

\begin{proof}
Событие $\{\delta_b>x\}$ при $x<W_b$ означает, что на интервале
$(\ln L_b-x,\ln L_b]$ длины $x$ нет ни одного события процесса; вероятность этого
равна $e^{-\lambda_b x}$. При $x\to W_b^-$ остаётся случай <<событий нет вовсе>>,
что и даёт атом в~\eqref{eq:delta}. Интегрируя хвост,
$\mathbb E[\delta_b]=\int_0^{W_b}e^{-\lambda_b x}\,dx=(1-e^{-\lambda_b W_b})/\lambda_b$,
откуда~\eqref{eq:Edelta}. Далее, $S_{\mathrm{full}}$ не случайна: по построению
$D_b+\delta_b=W_b$, поэтому $\kappa S_{\mathrm{full}}=\sum_b\lambda_b W_b=\mathbb E[M]$
--- это в точности уравнение моментов. Наконец,
$\kappa\,\mathbb E[\delta_b]/\ln b=\lambda_b\mathbb E[\delta_b]=1-e^{-\lambda_b W_b}$,
и суммирование по $b$ даёт~\eqref{eq:EBd}.
\end{proof}

\begin{remark}\label{rem:atom}
Отбрасываемый в~\eqref{eq:EBd} остаток $\sum_b e^{-\lambda_b W_b}$ есть ожидаемое
число оснований, у которых поиск не дал ни одной находки. В наших данных находки есть
у всех, а $\lambda_b W_b\approx N_b\ge7$, так что остаток не превосходит
$\sum_b e^{-N_b}<3\cdot10^{-3}$ против $B=15$ --- относительная поправка порядка
$0{,}02\%$. Именно её и показывает первая строка таблицы~\ref{tab:frontier}.
Подчеркнём, что \emph{без} этого усечения утверждение $\delta_b\sim\mathrm{Exp}
(\lambda_b)$ было бы неверным: у процесса нет событий левее $t_0$, поэтому
$\delta_b\le W_b$ по построению, и ссылка на отсутствие последействия здесь
недостаточна.
\end{remark}

Иными словами, условное правдоподобие стоит информации, эквивалентной одному событию
\emph{на каждое основание}, --- а формула~\eqref{eq:mle} вычитает одно событие
\emph{всего}. При $B=15$ разница составляет четырнадцать событий, то есть $7\%$:
\begin{equation}
\frac{M-1}{S}=1{,}9648\qquad\text{против}\qquad\frac{M-B}{S}=1{,}8381 .
\end{equation}
Прямая симуляция схемы~B подтверждает: оценка~\eqref{eq:MBS} имеет смещение $+0{,}02\%$,
тогда как $(M-1)/S$ --- смещение $+6{,}9\%$ (таблица~\ref{tab:calib}).

Существенно, что вывод предложения~\ref{prop:schemeB} опирается \emph{только} на
равенство~\eqref{eq:Edelta}, которое верно для \emph{любого} положения фронта:
в него входит лишь ширина окна $W_b$, а она сокращается вместе с
экспоненциально малым остатком (замечание~\ref{rem:atom}). Поэтому
оценка~\eqref{eq:MBS} не содержит свободного параметра и не требует знания того,
насколько далеко ушёл фронт. Таблица~\ref{tab:frontier} подтверждает это симуляцией
при семи различных правилах расположения фронта, включая случайные и различные по
основаниям.

\begin{table}[t]
\centering
\caption{Несмещённость двух объединённых оценок при различном расположении фронта.
Истинное значение $\eg=1{,}7811$; $15\,000$ реализаций на строку.
Правила заданы через ширину окна $W_b=\ln L_b-t_0$: <<$\mathbb E[\#]$ = наблюдённое>>
есть $W_b=N_b/\lambda_b$; <<$+s$ средних интервалов>> --- $W_b=(N_b+s)/\lambda_b$;
<<в $c$ раз дальше по $t$>> --- $W_b=c\,D_b$ с наблюдённым $D_b$; <<случайные>> ---
$W_b=(N_b/\lambda_b)\,U$ с независимой $U\sim\mathrm{Unif}(0{,}7;\,1{,}8)$, своей
у каждого основания и каждой реализации.
\encaptab{Unbiasedness of the two pooled estimators under different frontier
placements. True value $\eg=1.7811$; $15\,000$ replications per row. The rules are
stated through the window width $W_b=\ln L_b-t_0$: ``$\mathbb E[\#]$ = observed'' is
$W_b=N_b/\lambda_b$; ``$+s$ mean intervals'' is $W_b=(N_b+s)/\lambda_b$; ``$c$ times
further in $t$'' is $W_b=c\,D_b$ with the observed $D_b$; ``random'' is
$W_b=(N_b/\lambda_b)\,U$ with independent $U\sim\mathrm{Unif}(0.7,\,1.8)$, drawn
separately for each base and each replication.}}
\label{tab:frontier}
\small
\begin{tabular}{lrrrr}
\toprule
Правило для фронта & $\mathbb E[(M{-}B)/S]$ & откл. & $\mathbb E[(M{-}1)/S]$ & откл.\\
\midrule
$\mathbb E[\#\text{событий}]$ = наблюдённое & 1{,}7814 & $+0{,}02\%$ & 1{,}9043 & $+6{,}92\%$\\
фронт $+0{,}5$ среднего интервала & 1{,}7816 & $+0{,}03\%$ & 1{,}9001 & $+6{,}68\%$\\
фронт $+2$ средних интервала & 1{,}7820 & $+0{,}05\%$ & 1{,}8891 & $+6{,}07\%$\\
фронт $+5$ средних интервалов & 1{,}7820 & $+0{,}05\%$ & 1{,}8717 & $+5{,}09\%$\\
фронт в $1{,}5$ раза дальше по $t$ & 1{,}7825 & $+0{,}08\%$ & 1{,}8716 & $+5{,}08\%$\\
фронт вдвое дальше по $t$ & 1{,}7816 & $+0{,}03\%$ & 1{,}8475 & $+3{,}73\%$\\
фронты случайные, разные по основаниям & 1{,}7812 & $+0{,}01\%$ & 1{,}8788 & $+5{,}49\%$\\
\bottomrule
\end{tabular}
\end{table}

\paragraph{Единственное допущение.} Требуется лишь следующее.

\medskip
\noindent\textbf{(A1)} \emph{Положение фронта $L_b$ не выбиралось по результатам поиска:
вычислительный бюджет, а не находки, определяет, до какого $n$ доведено просеивание.}

\medskip
\noindent Если (A1) нарушено в крайней форме --- поиск по каждому основанию прекращался
сразу после находки, --- то $\delta_b=0$ и мы возвращаемся к схеме~A. Общий вид оценки
при $\delta_b=m/\lambda_b$ есть $\hat\kappa=(M-Bm)/S$, так что весь диапазон,
порождаемый неопределённостью в (A1), есть отрезок между $1{,}838$ ($m=1$, фронт
независим от находок) и $1{,}974$ ($m=0$, поиск останавливается на находке).\footnote{%
При $m=0$ моментное приравнивание даёт ровно $M/S=1{,}9739$, тогда как точная
несмещённая оценка схемы~A из следствия~\ref{cor:kappa} есть $(M-1)/S=1{,}9648$
(различие --- одно событие из $218$, то есть $0{,}5\%$; это обычная разница между
моментной оценкой и оценкой, исправленной по Йенсену). В~\S\ref{sec:pooled} и всюду
далее мы приводим $(M-1)/S$; здесь для однородности семейства $(M-Bm)/S$ выписан
предел моментной формы.}
Реальная практика распределённых вычислений явно ближе к первому: просеивание ведётся
сплошными интервалами, а находки объявляются по ходу. Существенно, что оба конца
отрезка совместимы с ЛПВ (\S\ref{sec:pooled}), так что вывод работы от (A1) не зависит;
от (A1) зависит лишь величина остаточного разрыва.

\begin{table}[t]
\centering
\caption{Калибровка пивота~\eqref{eq:pivot} при схеме~B. <<Сдвиг фронта>> --- на сколько
средних интервалов $\ln L_b$ отнесён дальше ожидаемого последнего события.
$20\,000$ реализаций на строку; точность Монте-Карло около $\pm0{,}005$
(биномиальная при наблюдаемом уровне $0{,}16$; разброс по пяти зёрнам
генератора --- $0{,}162$--$0{,}169$).
\encaptab{Calibration of pivot~\eqref{eq:pivot} under scheme~B. ``Frontier shift''
is how many mean intervals beyond the expected last event $\ln L_b$ is placed.
$20\,000$ replications per row; Monte Carlo accuracy about $\pm0.005$ (binomial at
the observed level $0.16$; the spread across five generator seeds is
$0.162$--$0.169$).}}
\label{tab:calib}
\small
\begin{tabular}{lrr}
\toprule
Схема & ошибка I рода (номинал $0{,}05$) & $\mathbb E[\tilde\kappa]/\kappa=c$\\
\midrule
B, сдвиг $0$ & 0{,}1616 & 1{,}0688\\
B, сдвиг $0{,}5$ & 0{,}1600 & 1{,}0664\\
B, сдвиг $1$ & 0{,}1558 & 1{,}0641\\
B, сдвиг $2$ & 0{,}1491 & 1{,}0595\\
A (контроль) & 0{,}0504 & 0{,}9997\\
\bottomrule
\end{tabular}
\end{table}

Два вывода из таблицы~\ref{tab:calib}. Во-первых, номинальный $5\%$ критерий при
схеме~B отвергает верную гипотезу в $16{,}2\%$ случаев: заявленная точность интервалов
завышена примерно втрое. Во-вторых, оценка~\eqref{eq:MBS} при схеме~B практически
несмещена, так что дело не в неустранимом искажении, а в неверно выбранной схеме.

Единственная прямая проверка возможна для $b=2$: там
$\delta_2=\ln(X_{\mathrm{ver}}/n_{\max})=0{,}056$ при среднем интервале
$\ln 2/\eg=0{,}389$, предсказанном ЛПВ.
Это одна реализация экспоненциальной величины с единичным средним, и такое значение
не редкость: $\Pr(\delta\lambda<0{,}143)=0{,}13$. (Если брать не предсказанный,
а оценённый средний интервал $1/\hat\lambda=0{,}358$, то $\delta\lambda=0{,}155$
и $\Pr=0{,}14$; вывод тот же.) Отметим, что несмещённость
оценки~\eqref{eq:MBS} от этой реализации не зависит --- она усредняет по всем
пятнадцати основаниям, --- а сама изменчивость $\delta_b$ учтена в доверительном
интервале~\eqref{eq:resB}, полученном инверсией симуляции схемы~B.

\subsection{Данные не выбирают схему за нас}\label{sec:noescape}

Подчеркнём, что $(M-1)/S$ и $(M-B)/S$ --- функции одних и тех же достаточных статистик
$(N_b,D_b)$. Никакая обработка имеющихся данных не скажет, какая из них верна: выбор
определяется тем, как данные собирались, и цена его составляет $7\%$ --- вдвое больше,
чем весь остаточный разрыв с ЛПВ.

Естественная попытка обойти выбор --- строить объединённую оценку из заведомо
несмещённых кусков.
Оценка по одному основанию $\tilde\kappa_b=(N_b-1)\ln b/D_b$ несмещена при схеме~A
(следствие~\ref{cor:kappa}) и несмещена при схеме~B с точностью до экспоненциально
малого члена. Действительно, при $N_b=n\ge2$ отношение $D_b/W_b$ распределено как
$\mathrm{Beta}(n,1)$, и тождество
$\mathbb E[1/\mathrm{Beta}(\alpha,\beta)]=(\alpha+\beta-1)/(\alpha-1)$ даёт
$\mathbb E[\tilde\lambda_b\mid N_b=n]=n/W_b$; при $n\in\{0,1\}$ оценка равна нулю.
Усреднение по $n\sim\mathrm{Pois}(\lambda_bW_b)$ поэтому даёт не $\lambda_b$, а
\[
  \mathbb E[\tilde\lambda_b]=\lambda_b-\frac{\Pr(N_b=1)}{W_b}
  =\lambda_b\Bigl(1-e^{-\lambda_bW_b}\Bigr),
\]
то есть тот же экспоненциально малый дефицит, что и в~\eqref{eq:EBd}: при
$\lambda_bW_b\approx N_b\ge7$ он не превосходит $0{,}1\%$. Симуляция подтверждает: по каждому из пятнадцати оснований среднее
$\tilde\kappa_b$ лежит в пределах $1{,}774$--$1{,}790$ при истинном $1{,}781$
(оценка имеет бесконечную дисперсию при $N_b=2$, поэтому выборочное среднее
сходится медленно; остаточный разброс по основаниям --- шум Монте-Карло).

Тем не менее взвешенное среднее
$$\kappa_\ast=\sum\nolimits_b(N_b-2)\tilde\kappa_b\Big/\sum\nolimits_b(N_b-2)$$
с обратно-дисперсионными весами даёт при схеме~A смещение $+0{,}05\%$,
а при схеме~B --- $+7{,}51\%$, то есть практически то же, что и $(M-1)/S$. Причина в том, что веса --- сами числа событий ---
при схеме~B случайны и положительно коррелируют с оценками; несмещённость каждого
слагаемого не переносится на отношение сумм.

Отсюда мы делаем узкий, но, как нам кажется, важный вывод: \emph{выигрыш в точности от
объединения оснований неотделим от предположения о схеме наблюдения}.
Единственный радикальный выход --- документировать фронты поиска $L_b$, как это делает
GIMPS для $b=2$. Мы обращаемся с этим к участникам соответствующих поисковых проектов.

\section{Результаты}\label{sec:results}

\subsection{Оценки по основаниям}

\begin{table}[t]
\centering
\caption{Оценки по основаниям. $\hat G_b=D_b/(N_b\ln b)$ --- концевая форма;
$\hat G_b^{\mathrm{OLS}}$ --- наклон МНК $\ln n_k$ на $k$, делённый на $\ln b$; обе
несмещены (предложение~\ref{prop:unbiased}). $\tilde\kappa_b=(N_b-1)\ln b/D_b$ ---
несмещённая оценка константы, ДИ точный из $\Gam(N_b,1)$. Последний столбец ---
вероятностное интегральное преобразование $\Gam(N_b,1)$, применённое к
$\eg D_b/\ln b$. ЛПВ предсказывает $G=\emg=0{,}5615$ и $\kappa=\eg=1{,}7811$.
\encaptab{Per-base estimates. $\hat G_b=D_b/(N_b\ln b)$ is the endpoint form;
$\hat G_b^{\mathrm{OLS}}$ is the OLS slope of $\ln n_k$ on $k$ divided by $\ln b$;
both are unbiased (Proposition~\ref{prop:unbiased}). $\tilde\kappa_b=(N_b-1)\ln b/D_b$
is the unbiased estimator of the constant, with an exact interval from $\Gam(N_b,1)$.
The last column is the probability integral transform of $\Gam(N_b,1)$ applied to
$\eg D_b/\ln b$. LPW predicts $G=\emg=0.5615$ and $\kappa=\eg=1.7811$.}}
\label{tab:bases}
\small
\begin{tabular}{rrrrrrrl r}
\toprule
$b$ & $N_b$ & $\hat G_b$ & $\hat G_b^{\mathrm{OLS}}$ & разн., \% &
$\hat\kappa_b$ & $\tilde\kappa_b$ & 95\% ДИ & PIT\\
\midrule
2 & 51 & 0{,}5102 & 0{,}5309 & $+4{,}1$ & 1{,}960 & 1{,}921 & [1{,}459; 2{,}533] & 0{,}266\\
3 & 22 & 0{,}6148 & 0{,}6081 & $-1{,}1$ & 1{,}626 & 1{,}552 & [1{,}019; 2{,}373] & 0{,}693\\
5 & 18 & 0{,}4802 & 0{,}4666 & $-2{,}8$ & 2{,}083 & 1{,}967 & [1{,}234; 3{,}149] & 0{,}285\\
6 & 16 & 0{,}5000 & 0{,}4904 & $-1{,}9$ & 2{,}000 & 1{,}875 & [1{,}143; 3{,}092] & 0{,}355\\
7 & 9 & 0{,}7104 & 0{,}7371 & $+3{,}8$ & 1{,}408 & 1{,}251 & [0{,}644; 2{,}466] & 0{,}801\\
10 & 10 & 0{,}6612 & 0{,}6746 & $+2{,}0$ & 1{,}512 & 1{,}361 & [0{,}725; 2{,}584] & 0{,}737\\
11 & 12 & 0{,}4034 & 0{,}3617 & $-10{,}3$ & 2{,}479 & 2{,}272 & [1{,}281; 4{,}066] & 0{,}162\\
12 & 13 & 0{,}3981 & 0{,}3796 & $-4{,}7$ & 2{,}512 & 2{,}319 & [1{,}337; 4{,}050] & 0{,}141\\
13 & 12 & 0{,}4098 & 0{,}3448 & $-15{,}9$ & 2{,}440 & 2{,}237 & [1{,}261; 4{,}002] & 0{,}174\\
14 & 10 & 0{,}5025 & 0{,}5386 & $+7{,}2$ & 1{,}990 & 1{,}791 & [0{,}954; 3{,}400] & 0{,}406\\
15 & 9 & 0{,}5035 & 0{,}5054 & $+0{,}4$ & 1{,}986 & 1{,}766 & [0{,}908; 3{,}479] & 0{,}417\\
17 & 11 & 0{,}4301 & 0{,}4353 & $+1{,}2$ & 2{,}325 & 2{,}114 & [1{,}161; 3{,}887] & 0{,}229\\
18 & 7 & 0{,}6604 & 0{,}4656 & $-29{,}5$ & 1{,}514 & 1{,}298 & [0{,}609; 2{,}825] & 0{,}714\\
19 & 10 & 0{,}3161 & 0{,}2971 & $-6{,}0$ & 3{,}164 & 2{,}847 & [1{,}517; 5{,}405] & 0{,}061\\
20 & 8 & 0{,}5300 & 0{,}5713 & $+7{,}8$ & 1{,}887 & 1{,}651 & [0{,}815; 3{,}402] & 0{,}483\\
\bottomrule
\end{tabular}
\end{table}

Ниже $\emg$ лежат $11$ оснований из $15$ --- одинаково при обеих формах оценки, так что
знак отклонения устойчив к её выбору, в отличие от порядка оснований между собой.
Все пятнадцать PIT-квантилей лежат внутри центральной $95\%$ области; крайние --- $b=19$
на квантиле $0{,}061$ и $b=7$ на $0{,}801$. Основание~13, <<удачливое>>, находится на
$0{,}174$, что ничем не примечательно для одного из пятнадцати.

\subsection{Объединённая оценка}\label{sec:pooled}

При схеме~A формулы~\eqref{eq:mle}--\eqref{eq:pivot} дают
\begin{equation}\label{eq:resA}
\tilde\kappa=1{,}9648,\qquad 95\%\ \text{ДИ}=[1{,}7205;\,2{,}2444],\qquad p=0{,}142 .
\end{equation}
Точечная оценка на $10\%$ выше $\eg$, и именно это превышение прежде интерпретировалось
как след опущенного эвристикой фактора.

Оценка~\eqref{eq:MBS} для схемы~B даёт (интервал и $p$-значение получены инверсией
симуляции схемы~B: для сетки значений $\kappa$ симулируется схема~B с наблюдённым
расположением фронта и отыскиваются те $\kappa$, при которых наблюдённая статистика
$(M-B)/S$ попадает на квантили $0{,}025$ и $0{,}975$ симулированного распределения;
$20\,000$ реализаций на точку)
\begin{equation}\label{eq:resB}
\hat\kappa_{\mathrm B}=\frac{M-B}{S}=1{,}8381,\qquad
95\%\ \text{ДИ}=[1{,}613;\,2{,}130],\qquad p=0{,}639 .
\end{equation}
Превышение сокращается с $10{,}3\%$ до $3{,}2\%$, а $p$-значение --- с $0{,}14$ до
$0{,}64$. Мы полагаем, что схема~B ближе к способу получения данных, и потому
считаем~\eqref{eq:resB} основным результатом, а~\eqref{eq:resA} --- верхней границей:
она отвечает предположению $m=0$, то есть остановке поиска ровно на находке.
Истина лежит между ними, и без публикации $L_b$ сузить этот промежуток нельзя.
Обе оценки согласуются с ЛПВ; различаются они тем, требуют ли данные объяснения
для остаточного превышения. Мы заключаем, что не требуют. Картина в целом показана
на рис.~\ref{fig:kappa}: доверительный интервал каждого основания накрывает
и~$\eg$, и обе объединённые оценки.

\begin{figure}[t]
\centering
\includegraphics[width=.85\textwidth]{fig_kappa.pdf}
\caption{Оценки $\tilde\kappa_b$ по основаниям с точными $95\%$ доверительными
интервалами из $\Gam(N_b,1)$; площадь маркера пропорциональна $N_b$. Штриховая линия ---
значение ЛПВ $\eg$; синяя --- объединённая оценка при схеме~A; зелёная с полосой ---
калиброванная на схему~B оценка и её интервал. Каждое основание по отдельности
совместимо с универсальной константой.
\encapfig{Per-base estimates $\tilde\kappa_b$ with exact $95\%$ confidence
intervals from $\Gam(N_b,1)$; marker area is proportional to $N_b$. Dashed line: the
LPW value $\eg$; blue dash-dotted: the pooled estimate under scheme~A; green with a
hatched band: the estimate calibrated to scheme~B and its interval. Each base
separately is compatible with a universal constant.}}
\label{fig:kappa}
\end{figure}

\subsection{Однородность и её мощность}\label{sec:homog}

Пусть $H_0$ --- модель с единой $\kappa$, а $H_1$ --- модель со свободными $\kappa_b$.
Из~\eqref{eq:loglik}, поскольку $\sum_b\hat\kappa_bD_b/\ln b=M=\hat\kappa S$, статистика
отношения правдоподобий сводится к
\begin{equation}
\Lambda=2\sum_b N_b\ln\frac{\hat\kappa_b}{\hat\kappa}\ \overset{H_0}{\sim}\ \chi^2_{B-1}.
\end{equation}
Мы получаем $\Lambda=7{,}40$ при $14$ степенях свободы, $p=0{,}92$.

Однако неотвержение информативно ровно настолько, насколько мощен критерий, и здесь
следует применить к себе то же требование, что и к литературе. Мы симулировали
альтернативы, в которых $\kappa_b$ логнормально рассеяны вокруг $\eg$ с коэффициентом
вариации $cv$:

\begin{center}
\small
\begin{tabular}{lrrrrrrrrr}
\toprule
$cv$ & 0{,}05 & 0{,}10 & 0{,}15 & 0{,}20 & 0{,}25 & 0{,}30 & 0{,}35 & 0{,}40 & 0{,}50\\
мощность & 0{,}061 & 0{,}120 & 0{,}212 & 0{,}363 & 0{,}535 & 0{,}693 & 0{,}817 & 0{,}893 & 0{,}965\\
\bottomrule
\end{tabular}
\end{center}

Мощности $80\%$ критерий достигает лишь при $cv\approx0{,}34$ ($5000$ реализаций
на точку сетки, точность по мощности около $\pm0{,}01$). Следовательно,
$p=0{,}92$ означает: данные не требуют базозависимых частот, но и не исключают
разброса $\kappa_b$ величиной до трети от значения константы. Утверждение
<<удачливых оснований не существует>> из имеющихся данных не следует; следует более
слабое --- <<их существование данными не поддерживается, а различия менее $\pm34\%$
принципиально неразличимы>>.

\subsection{Устойчивость}\label{sec:robust}

\paragraph{Нижнее усечение.} Мы варьировали нижнюю точку усечения, чтобы прощупать
асимптотический режим (эвристика~\eqref{eq:LPW} асимптотична, тогда как модель
постоянной интенсивности стартует от $n=2$). Результаты в таблице~\ref{tab:trunc}.
Оценка схемы~A растёт монотонно с $1{,}965$ до $2{,}099$; это систематический тренд,
а не шум. Оценка схемы~B тренда не показывает и колеблется вокруг $\eg$
в пределах $1{,}79$--$1{,}87$ (при более слабом критерии выбытия $N_b\ge1$ ---
в пределах $1{,}76$--$1{,}87$, по-прежнему без тренда). Причина видна
из~\eqref{eq:MBS}: усечение уменьшает $M$, но не $B$, тогда как схема~A вычитает
единицу независимо от объёма выборки. Мы считаем это независимым подтверждением
диагноза~\S\ref{sec:scheme}: тренд был свойством оценки, а не данных.
Обе кривые приведены на рис.~\ref{fig:trunc}.

\begin{table}[t]
\centering
\caption{Устойчивость к нижнему усечению. Основание остаётся в выборке, если после
усечения у него сохраняется не менее \emph{двух} событий ($N_b\ge2$); при $n>10^4$ этому
критерию перестаёт удовлетворять основание~19, откуда $B=14$. Критерий существен:
при более слабом требовании $N_b\ge1$ основание~19 сохраняется, и последняя строка
принимает вид $B=15$, $M=90$, $(M-1)/S=2{,}0848$, $(M-B)/S=1{,}7569$, $p=0{,}128$.
\encaptab{Robustness to lower truncation. A base is retained if at least
\emph{two} events survive the truncation ($N_b\ge2$); at $n>10^4$ base~19 no longer
meets this criterion, whence $B=14$. The criterion matters: under the weaker
requirement $N_b\ge1$ base~19 is retained and the last row becomes $B=15$, $M=90$,
$(M-1)/S=2.0848$, $(M-B)/S=1.7569$, $p=0.128$.}}
\label{tab:trunc}
\small
\begin{tabular}{lrrrrr}
\toprule
усечение & $B$ & $M$ & $(M-1)/S$ & $(M-B)/S$ & $p$ (схема A)\\
\midrule
нет & 15 & 218 & 1{,}9648 & 1{,}8381 & 0{,}142\\
$n>10$ & 15 & 192 & 2{,}0192 & 1{,}8712 & 0{,}080\\
$n>10^2$ & 15 & 156 & 2{,}0116 & 1{,}8299 & 0{,}124\\
$n>10^3$ & 15 & 125 & 2{,}0551 & 1{,}8231 & 0{,}106\\
$n>10^4$ & 14 & 89 & 2{,}0990 & 1{,}7889 & 0{,}116\\
\bottomrule
\end{tabular}
\end{table}

\begin{figure}[t]
\centering
\includegraphics[width=.75\textwidth]{fig_truncation.pdf}
\caption{Объединённая оценка при нижнем усечении: $(M-1)/S$ (схема~A) и $(M-B)/S$
(схема~B). Тренд оценки схемы~A при переходе к схеме~B исчезает.
\encapfig{The pooled estimate under lower truncation: $(M-1)/S$ (scheme~A) and
$(M-B)/S$ (scheme~B). The trend of the scheme~A estimate disappears under scheme~B.}}
\label{fig:trunc}
\end{figure}

\paragraph{Исключение по одному основанию.} Анализ с исключением одного основания даёт
$\tilde\kappa$ в пределах $[1{,}930;\,2{,}012]$ и $p$ от $0{,}086$ до $0{,}237$;
ни одно основание не определяет вывод, а наиболее влиятельное, $b=3$, при удалении
сдвигает $p$ лишь до $0{,}086$.

\paragraph{Статус данных.} Ограничение основания~2 пятьюдесятью дважды проверенными
показателями переводит результат с $\tilde\kappa=1{,}965$, $p=0{,}142$ на
$\tilde\kappa=1{,}961$, $p=0{,}151$. Исключение всех индексов $n>10^6$, где статус
элементов --- PRP, даёт $\tilde\kappa=1{,}951$ при $M=182$. В наихудшем сценарии, когда
последний известный PRP какого-либо основания оказывается составным, наибольший сдвиг
даёт $b=17$: $\tilde\kappa=1{,}977$. Все варианты лежат внутри интервала~\eqref{eq:resA},
поэтому возможные псевдопростые на вывод не влияют.

\section{Статистическая мощность}\label{sec:power}

Согласно~\eqref{eq:pivot} мощность против истинной константы $\kappa=r\eg$ на уровне
$\alpha$ доступна в замкнутой форме (таблица~\ref{tab:power}).

\begin{table}[t]
\centering
\caption{Мощность на уровне $\alpha=0{,}05$ для отвержения $\kappa=\eg$, когда истинная
константа превышает его в указанное число раз.
\encaptab{Power at level $\alpha=0.05$ for rejecting $\kappa=\eg$ when the true
constant exceeds it by the stated factor.}}
\label{tab:power}
\small
\begin{tabular}{lrr}
\toprule
$\kappa/\eg$ & только Мерсенн ($M=51$) & объединение ($M=218$)\\
\midrule
1{,}05 & 0{,}060 & 0{,}106\\
1{,}10 & 0{,}095 & 0{,}277\\
1{,}15 & 0{,}152 & 0{,}523\\
1{,}20 & 0{,}229 & 0{,}755\\
1{,}30 & 0{,}427 & 0{,}972\\
1{,}50 & 0{,}802 & 1{,}000\\
\bottomrule
\end{tabular}
\end{table}

При мощности $80\%$ одна последовательность Мерсенна обнаруживает только отклонения
в $49{,}8\%$ и более; объединение пятнадцати оснований снижает этот порог до $21{,}2\%$.
Достижение мощности $80\%$ против отклонения в $10\%$ требует $M\approx875$ событий,
против $5\%$ --- $M\approx3319$.

Эти рубежи вычислительно недостижимы, и стоит сказать насколько. Поскольку $M$ растёт
логарифмически по фронту поиска, удвоение фронтов \emph{по всем пятнадцати основаниям}
даёт около $9{,}6$ дополнительных событий на всю объединённую выборку. Чтобы набрать
недостающие $657$ событий, нужно $657/9{,}6\approx68$ удвоений, то есть рост фронтов
в $2^{68{,}2}\approx3\cdot10^{20}$ раз. Речь идёт
не о десятилетиях, а о величине, лежащей вне досягаемости любых мыслимых вычислений.
Обе кривые мощности как функции истинного отклонения показаны на
рис.~\ref{fig:power}.

\begin{figure}[t]
\centering
\includegraphics[width=.72\textwidth]{fig_power.pdf}
\caption{Мощность отвержения $\kappa=\eg$ на уровне $\alpha=0{,}05$. Горизонтальная
линия отмечает мощность $80\%$.
\encapfig{Power of rejecting $\kappa=\eg$ at level $\alpha=0.05$. The horizontal
line marks $80\%$ power.}}
\label{fig:power}
\end{figure}

Отсюда правило чтения для этой литературы. Сообщение о том, что данные согласуются с ЛПВ
с точностью до нескольких процентов, не есть свидетельство в пользу ЛПВ на уровне
нескольких процентов; оно совместимо с константой, ошибочной в полтора раза.
Объединённый анализ устанавливает, однако, настоящую границу: универсальная константа,
если она существует, лежит примерно в $[1{,}61;\,2{,}13]$, и ни одно основание от неё
различимо не отклоняется.

\section{Та же ошибка в ином обличье}\label{sec:sameerror}

Счётная функция Мерсенна провоцирует родственную ошибку, которую мы приводим как
разобранный пример. Регрессия $\pi_2(t)$ на~$t$ в точках событий по $50$ дважды
проверенным показателям (нумерация $0,\dots,N-1$: наименьший показатель есть начало
отсчёта, а не событие) даёт наклон $2{,}6391$ и свободный член $-2{,}9627$ при
стандартной ошибке МНК $0{,}4469$, на вид значимый и наводящий на мысль о вторичном
члене в~\eqref{eq:LPW}. Ничего подобного.

Складываются два дефекта. Остатки кумулятивного счёта сильно автокоррелированы,
здесь $\rho_1=0{,}825$; ковариационная оценка Ньюи --- Уэста с восемью лагами
увеличивает стандартную ошибку до $0{,}7315$. Существеннее то, что регрессия использует
только точки событий и тем самым отбрасывает пустой интервал между последним показателем
и фронтом поиска. Корректная нулевая модель удерживает интервал $[t_0,T]$ фиксированным,
где $T=\ln X_{\mathrm{ver}}$, и обусловливается ровно $N-1=49$ событиями строго правее
$t_0$, что по предложению~\ref{prop:gamma} означает $49$ независимых равномерных точек
на $(t_0,T]$. По $20\,000$ таким реализациям свободный член имеет медиану $-1{,}498$ и
$5$--$95\%$ диапазон $[-5{,}99;\,2{,}20]$, что даёт $p=0{,}562$ для наблюдённого значения.
Вторичный член есть артефакт обусловливания.

\subsection{Три следствия одного и того же обусловливания}

\paragraph{(i) Бутстрэп не спасает.} Естественная реакция на подозрительно малую
стандартную ошибку МНК --- бутстрэп по точкам событий. На данных основания~2
бутстрэп по $20\,000$ ресемплам даёт медиану $-2{,}949$, стандартное отклонение
$0{,}387$ и $95\%$ интервал $[-3{,}822;\,-2{,}300]$, целиком лежащий ниже нуля; доля отрицательных значений
равна $100\%$. Это не подтверждение, а повторение той же ошибки: ресемплинг наблюдённых
точек воспроизводит изменчивость внутри реализованной конфигурации и не может
восстановить изменчивость самого процесса, породившего эту конфигурацию. Корректное
нулевое распределение примерно в шесть с половиной раз шире бутстрэпного.

\paragraph{(ii) Мнимая недодисперсия.} Сравнение наблюдённого $\max_i|r_i|=6{,}347$,
где $r_i=i-\hat\lambda(t_i-t_0)$, с симуляцией \emph{необусловленного} пуассоновского
процесса наводит на вывод, что счётная функция <<глаже пуассоновской>>. Против
корректного нулевого распределения --- с обусловливанием на $N-1$ событий и оценкой
$\lambda$ по последнему из них --- медиана равна $5{,}134$, а двусторонний уровень
значимости составляет $0{,}539$. Наблюдение совершенно рядовое.

\paragraph{(iii) <<Серединный дефицит>>.} Оценивание $\lambda$ по последнему событию
обращает остаток $r$ в нуль на обоих концах ($r(t_0)=r(t_N)=0$ при правильной
нумерации), превращая его траекторию в аналог броуновского моста, для которого
отрицательная экскурсия в середине диапазона возникает сама собой. Наблюдённая средняя
экскурсия на средней половине диапазона равна $-3{,}13$ при медиане симуляций
$-0{,}001$ и $p=0{,}280$; доля реализаций с отрицательной средней экскурсией
составляет $0{,}500$, то есть систематического дефицита нет. Сама траектория
остаточного процесса вместе с полосой симуляций приведена на
рис.~\ref{fig:bridge}.

\begin{figure}[t]
\centering
\includegraphics[width=.78\textwidth]{fig_resid_bridge.pdf}
\caption{Остаточный процесс $r(t)$ для основания $b=2$ вместе с полосой $5$--$95\%$
квантилей $20\,000$ симуляций и медианой симуляций. Траектория лежит внутри полосы.
\encapfig{The residual process $r(t)$ for base $b=2$ together with the
$5$--$95\%$ quantile band of $20\,000$ simulations and the simulation median. The
trajectory stays inside the band.}}
\label{fig:bridge}
\end{figure}

\subsection{Основание 2 на известном фронте}

Основание~2 --- единственное, где фронт опубликован, и потому единственное, где возможен
анализ по схеме~B без всяких предположений. С $T=\ln X_{\mathrm{ver}}$ получаем
$\hat\lambda=(N-1)/(T-t_0)=2{,}796$ против $\eg/\ln2=2{,}570$, то есть
$\hat\kappa_2=1{,}938$ при точном $95\%$ интервале Гарвуда $[1{,}434;\,2{,}562]$ и
пуассоновском $p=0{,}59$ (удвоенный хвост; при mid-$p$-конвенции $0{,}54$ --- конвенцию
здесь необходимо оговаривать, поскольку различие ощутимо). Условная оценка \emph{того же
среза данных} --- те же $50$ дважды проверенных показателей --- даёт
$\tilde\kappa_2=1{,}905$; различие невелико, поскольку последний показатель этого среза
$77\,232\,917$ случайно оказался очень близко к фронту $81\,648\,221$. (Для сравнения,
условная оценка по всем $52$ известным показателям, приведённая в
таблице~\ref{tab:bases}, равна $1{,}921$; сопоставлять её с $\hat\kappa_2$ нельзя ---
это другая выборка.) Это совпадение
одной реализации, а не свойство метода: смещение~\S\ref{sec:scheme} относится к среднему
по положению последнего события относительно фронта.

\subsection{Теоретический вторичный член имеет противоположный знак}

Существует независимый довод против содержательной трактовки свободного члена $C_1$.
Если вторичный член в~\eqref{eq:LPW} происходит из постоянной Мертенса
$\mathfrak M=0{,}26150\ldots$, то естественная его величина есть
$(\eg/\ln2)\mathfrak M=+0{,}672$: знак положительный, тогда как наблюдается
$-2{,}963$. Ближайшие теоретические кандидаты не просто не воспроизводят наблюдаемую
величину, а предсказывают противоположный знак. Это независимо от предыдущего
подтверждает, что $C_1$ --- параметр аппроксимации, а не структуры.

\section{Структурные поправки}\label{sec:struct}

Рассмотрим более общую параметризацию интенсивности $\lambda_b=(\kappa/\ln b)\,C_b$
и четыре модели множителя:
\begin{align*}
\text{A:}&\quad C_b=1;\\
\text{B:}&\quad C_b=\prod_{p\mid b-1,\ p\ne2}\Bigl(\tfrac{p}{p-1}\Bigr)^{\alpha}
                 \prod_{p\mid b+1,\ p\ne2}\Bigl(\tfrac{p}{p-2}\Bigr)^{\beta};\\
\text{C:}&\quad C_b=\prod_{p\mid b^2-1}\Bigl(\tfrac{p}{p-1}\Bigr)^{\delta};\\
\text{D:}&\quad C_b=\exp\bigl(\theta_1\omega(b-1)+\theta_2\omega(b+1)+\theta_3\ln b\bigr).
\end{align*}
Показатель в модели~C обозначен через $\delta$, а не $\gamma$, чтобы не смешивать его
с постоянной Эйлера. Результаты в таблице~\ref{tab:struct}. Ни одна параметризация не
даёт значимого улучшения; байесовские факторы по BIC при $n=M=218$ подавляюще
поддерживают базовую модель.

\begin{table}[t]
\centering
\caption{Сравнение структурных моделей. Правдоподобие --- условное~\eqref{eq:loglik}
с множителем $C_b$; LR и $p$ даны относительно модели~A; BIC вычислен при $n=M=218$.
\encaptab{Comparison of structural models. The likelihood is the
conditional one~\eqref{eq:loglik} with factor $C_b$; LR and $p$ are relative to
model~A; BIC is computed with $n=M=218$.}}
\label{tab:struct}
\small
\begin{tabular}{lrrrrrrr}
\toprule
Модель & параметров & $\ln L$ & LR & $p$ & AIC & BIC & $P(\text{модель}\mid\text{данные})$\\
\midrule
A & 1 & $-176{,}129$ & --- & --- & 354{,}26 & 357{,}64 & 0{,}893\\
B & 3 & $-175{,}494$ & 1{,}27 & 0{,}530 & 356{,}99 & 367{,}14 & 0{,}008\\
C & 2 & $-175{,}647$ & 0{,}96 & 0{,}326 & 355{,}29 & 362{,}06 & 0{,}098\\
D & 4 & $-174{,}925$ & 2{,}41 & 0{,}492 & 357{,}85 & 371{,}39 & 0{,}001\\
\bottomrule
\end{tabular}
\end{table}

\begin{remark}[Чувствительность BIC]
Выбор эффективного объёма выборки для BIC неочевиден: правдоподобие~\eqref{eq:loglik}
имеет $B=15$ независимых слагаемых, но информация растёт как $M=218$. При $n=15$
апостериорные вероятности составляют $0{,}624$ / $0{,}079$ / $0{,}261$ / $0{,}036$,
при $n=218$ --- $0{,}893$ / $0{,}008$ / $0{,}098$ / $0{,}001$. Вывод в пользу модели~A
при обоих выборах, но при $n=218$ он существенно сильнее. Мы отчитываемся по $n=M$
как соответствующему числу независимых событий.
\end{remark}

\paragraph{Фиксированная арифметическая поправка.} Отдельно рассмотрена модель с
\emph{фиксированным} (не оцениваемым) множителем
$C^{\mathrm{fix}}_b=\prod_{p\mid b-1,p\ne2}\frac{p}{p-1}\prod_{p\mid b+1,p\ne2}\frac{p}{p-2}$,
для которого $\overline{C_b}=2{,}46$. Она даёт $\hat\kappa=0{,}815$. Подчеркнём, что
сравнивать это значение с $\eg$ бессмысленно: в перепараметризованной модели $\kappa$
имеет иной масштаб, и такой <<критерий согласия>> проверял бы лишь равенство среднего
$C_b$ единице. Поскольку модели A и $C^{\mathrm{fix}}$ ненестированы и имеют одинаковое
число свободных параметров, корректное сравнение --- прямо по правдоподобию:
$\ln L=-193{,}55$ против $-176{,}13$, разность $-17{,}42$. Фиксированная поправка
решительно отвергается.

\paragraph{Логарифмические поправки.} Модель $\lambda_b(t)=(\kappa/\ln b)(1+c_1/t+c_2/t^2)$
оценивается полным правдоподобием точечного процесса
$\sum_i\ln\lambda(t_i)-\int\lambda$ (правдоподобие одного лишь счёта здесь
недопустимо: оно отбрасывает положения событий, которыми только и определяются
$c_1,c_2$). Получаем $\hat\kappa=2{,}049$, $\hat c_1=+0{,}108$, $\hat c_2=-0{,}826$,
$\mathrm{LR}=3{,}53$ при двух степенях свободы, $p=0{,}171$. Незначимо.

\section{Заключение}

\begin{enumerate}
\item Употребительные оценки нормированной частоты обобщённых репьюнитных простых
\emph{не} смещены схемой подгонки: и концевая форма, и наклон МНК несмещены точно.
Смещена лишь обратная величина, и это следствие выпуклости. Ходовое объяснение
<<медленной сходимости снизу>> через цензурирование несостоятельно.

\item Настоящий дефект употребительной практики --- в том, что различные формы оценки
расходятся до $29{,}5\%$, а точечные значения публикуются без доверительных интервалов
при том, что стандартные ошибки МНК в этой задаче занижены примерно в $1{,}6$ раза.
Классификация оснований на <<удачливые>> и <<неудачливые>> переворачивается от смены
формы оценки на неизменных данных.

\item Объединение $B=15$ оснований даёт $M=218$ событий и точный пивот
$\kappa S\sim\Gam(M,1)$, но этот пивот откалиброван под схему наблюдения <<до $N$-го
события>>, тогда как данные получены при фиксированном фронте поиска. При реальной
схеме номинальный $5\%$ критерий отвергает в $16{,}2\%$ случаев, а правильной
объединённой оценкой оказывается $(M-B)/S$: условное правдоподобие стоит
одного события на каждое основание, а не одного события в сумме.

\item Это даёт $\kappa=1{,}838$ при $95\%$ интервале $[1{,}61;\,2{,}13]$ и
$p=0{,}639$ против $\eg$. Сообщавшееся десятипроцентное превышение над ЛПВ, включая
наше собственное в более ранней редакции этой работы, есть артефакт схемы наблюдения.
Переход к схеме~B независимо устраняет и систематический тренд объединённой оценки
при усечении по $n_{\min}$.

\item Выбор между $(M-1)/S$ и $(M-B)/S$ не решается данными: обе оценки суть функции
одних и тех же достаточных статистик и различаются на $7\%$. Практическое следствие ---
необходимость публиковать фронты поиска $L_b$, как это делает GIMPS для $b=2$.
До тех пор любой объединённый анализ этих последовательностей несёт неустранимую
неопределённость порядка $7\%$ --- вдвое больше остаточного разрыва с ЛПВ.

\item Гипотеза ЛПВ имеющимися данными не отвергается, а базозависимые поправки не
требуются. Но и то и другое --- утверждения о низкой чувствительности критериев:
критерий однородности различает лишь разброс $\kappa_b$ порядка $\pm34\%$, а
объединённый критерий --- отклонения константы не менее $21\%$. Достижение
чувствительности в $10\%$ потребовало бы роста фронтов поиска примерно в $3\cdot10^{20}$ раз.
\end{enumerate}


\appendix
\section{Верификация исходных данных}\label{app:data}

Все двадцать последовательностей (пятнадцать основного набора и пять из проверки
границы отбора) сверены с записями OEIS 27 августа 2026~г.
Поскольку в объединённый анализ входят только тройки $(N_b,n_{\min},n_{\max})$,
в таблице~\ref{tab:verify} приведены именно они; полные списки индексов содержатся
в скрипте \texttt{analysis.py}.

\begin{table}[h]
\centering
\caption{Сверка с OEIS. Столбец <<последнее пополнение>> указывает происхождение
наибольшего известного индекса.
\encaptab{Cross-check against OEIS. The column ``last addition'' gives the
provenance of the largest known index.}}
\label{tab:verify}
\small
\begin{tabular}{rlrrrl}
\toprule
$b$ & OEIS & простых & $n_{\min}$ & $n_{\max}$ & последнее пополнение\\
\midrule
2 & A000043 & 52 & 2 & 136\,279\,841 & GIMPS, окт.\ 2024\\
3 & A028491 & 23 & 3 & 8\,530\,117 & P.~Bourdelais\\
5 & A004061 & 19 & 3 & 3\,300\,593 & P.~Bourdelais\\
6 & A004062 & 17 & 2 & 3\,360\,347 & P.~Bourdelais\\
7 & A004063 & 10 & 5 & 1\,264\,699 & P.~Bourdelais\\
10 & A004023 & 11 & 2 & 8\,177\,207 & Propper и Batalov, май 2021\\
11 & A005808 & 13 & 17 & 1\,868\,983 & P.~Bourdelais\\
12 & A004064 & 14 & 2 & 769\,543 & P.~Bourdelais, дек.\ 2014\\
13 & A016054 & 13 & 5 & 1\,503\,503 & P.~Bourdelais\\
14 & A006032 & 11 & 3 & 1\,724\,417 & P.~Bourdelais, 13.07.2026\\
15 & A006033 & 10 & 3 & 639\,833 & P.~Bourdelais, 22.04.2019\\
17 & A006034 & 12 & 3 & 1\,990\,523 & P.~Bourdelais, 03.08.2020\\
18 & A133857 & 8 & 2 & 1\,270\,141 & P.~Bourdelais\\
19 & A006035 & 11 & 19 & 209\,359 & P.~Bourdelais, 27.08.2010\\
20 & A127995 & 9 & 3 & 984\,349 & P.~Bourdelais\\
\bottomrule
\end{tabular}
\end{table}

Отдельно предупредим читателя, сверяющего данные по вторичным источникам: сводные
таблицы (страница <<Repunit primes>> в OEIS Wiki, статья <<List of repunit primes>>
в Википедии) обновляются с запаздыванием и на момент написания работы отставали
от самих записей OEIS на один член у оснований $b=14,15,17,18,20$. Авторитетным
источником следует считать записи OEIS.

\subsection*{Независимый пересчёт префиксов}

Сверка с OEIS подтверждает, что перечисленные индексы дают простое, но не
гарантирует \emph{полноты}: пропущенный член сдвинул бы $N_b$, а с ним и обе
объединённые оценки. Этот пробел мы закрыли прямым пересчётом. Для каждого из
двадцати оснований $2\le b\le26$, не являющихся точными степенями, независимой
реализацией\footnote{Гибридный CPU/GPU поисковик: сегментное сито по малым
делителям, trial factoring в арифметике Монтгомери, PRP через GMP и GWNUM/IBDWT.
Исходный код и протокол прогонов --- в том же репозитории, скрипты
\texttt{verify\_sequences.sh} и \texttt{compare\_verify.py}.} проверены
\emph{все} индексы $n<10^4$ подряд, а не только перечисленные в OEIS. Каждый
отрицательный вердикт пересчитан вторым, независимым бэкендом (сплошной
double-check, как в GIMPS), чтобы исключить потерянную находку.

Результат: $153$ члена подтверждены, лишних нет, пропущенных нет, расхождений
между бэкендами нет --- по всем двадцати основаниям. Покрытие по основаниям
неравномерно (от одного члена при $b=18$ до двадцати двух при $b=2$): на
основной набор $b\le20$ приходится $128$ подтверждённых индексов из $233$
известных. Выше $10^4$ полный пересчёт вычислительно недоступен, и там мы
по-прежнему опираемся на OEIS.

Отметим, что этот пересчёт не был формальностью: в первом прогоне независимая
реализация объявила $R_{3181}(23)$ составным, разойдясь с A204940. Разбор показал
ошибку в самой реализации --- всплеск ошибки округления FFT на последней итерации,
после последней периодической проверки, --- а не в данных OEIS. После исправления
расхождений не осталось. Мы приводим этот эпизод как иллюстрацию к тезису
\S\ref{sec:scheme}: ложное <<составное>> невидимо без повторного счёта, и именно
поэтому фронты и протоколы проверки должны публиковаться.

\section*{Доступность данных и кода}
Все использованные последовательности приведены в таблице~\ref{tab:data}.

\emph{Данные.} Скрипт \texttt{fetch\_data.sh} загружает b-файлы OEIS для всех
пятнадцати последовательностей основного набора, \texttt{fetch\_extra.sh} --- для
пяти оснований $21\le b\le26$, использованных в проверке границы отбора
(\S\ref{sec:bases}). Загруженные файлы включены в репозиторий, так что анализ
воспроизводится без обращения к сети.

\emph{Анализ.} Один самодостаточный скрипт \texttt{analysis.py} (numpy, scipy,
matplotlib); каждое число настоящей работы помечено в его выводе
(\texttt{results.txt}) тегом соответствующей таблицы или раздела и воспроизводится
одним запуском. Генератор случайных чисел инициализирован фиксированным зерном,
поэтому все величины Монте-Карло воспроизводятся побитово. Скрипт также строит
все четыре рисунка. Полный прогон занимает несколько минут.

\emph{Независимый пересчёт последовательностей.} \texttt{verify\_sequences.sh}
запускает поисковик по всем индексам $n<10^4$ для каждого основания,
\texttt{compare\_verify.py} сверяет результат с b-файлами OEIS и возвращает
ненулевой код при любом расхождении (приложение~\ref{app:data}).

\emph{Оговорка о величинах Монте-Карло.} Концы интервала~\eqref{eq:resB} зависят от
процедуры инверсии, а таблицы~\ref{tab:mc}, \ref{tab:frontier} и~\ref{tab:calib}~---
от правил расположения фронта. И то и другое описано в тексте и реализовано в
\texttt{analysis.py}; все приведённые в работе значения --- в точности вывод этого
скрипта при зерне, зафиксированном в нём. Точность симуляций: для уровня
значимости в таблице~\ref{tab:calib} около $\pm0{,}003$, для средних в
таблице~\ref{tab:frontier} --- около $\pm0{,}3\%$.

Репозиторий: \textcolor{red}{[ВСТАВИТЬ URL ПЕРЕД ПОДАЧЕЙ]} --- каталог
\texttt{paper/}. Дерево файлов и порядок запуска описаны в \texttt{paper/README.md}.

% п. 4.14 ГОСТ Р 7.0.7-2021: перечень затекстовых библиографических ссылок
% предваряют словами «Список источников»; babel-russian по умолчанию печатает
% «Список литературы», поэтому \refname переопределяем.
\renewcommand{\refname}{Список источников}
\begin{thebibliography}{9}

\bibitem{wagstaff1983} Wagstaff S.~S., Jr. Divisors of Mersenne numbers~//
Mathematics of Computation. 1983. Vol.~40, №~161. P.~385--397.
DOI: 10.1090/S0025-5718-1983-0679454-X.

\bibitem{pomerance1981} Pomerance C. On the distribution of pseudoprimes~//
Mathematics of Computation. 1981. Vol.~37, №~156. P.~587--593.
DOI: 10.1090/S0025-5718-1981-0628717-0.

\bibitem{t5k} Heuristics: deriving the Wagstaff Mersenne conjecture
[Электронный ресурс] // The PrimePages: prime number research and records.
URL: \url{https://t5k.org/mersenne/heuristic.html} (дата обращения: 27.08.2026).

\bibitem{dubner1993} Dubner H. Generalized repunit primes~//
Mathematics of Computation. 1993. Vol.~61, №~204. P.~927--930.
DOI: 10.1090/S0025-5718-1993-1185243-9.

\bibitem{bourdelais2009} Bourdelais P. A generalized repunit conjecture
[Электронный ресурс]~// NMBRTHRY mailing list. 25 June 2009.
URL: \url{https://listserv.nodak.edu/cgi-bin/wa.exe?A2=NMBRTHRY;417ab0d6.0906}
(дата обращения: 25.08.2026).

\bibitem{oeis} The On-Line Encyclopedia of Integer Sequences
[Электронный ресурс] / OEIS Foundation Inc.
URL: \url{https://oeis.org} (дата обращения: 27.08.2026).

\bibitem{wagstaff2025} Wagstaff S.~S., Jr. Two conjectures about Mersenne primes
[Электронный ресурс]~// Journal of Integer Sequences. 2025. Vol.~28.
Article 25.7.2.
URL: \url{https://cs.uwaterloo.ca/journals/JIS/VOL28/Wagstaff/wag3.html}
(дата обращения: 25.08.2026).

\bibitem{dominguez2026} Dominguez J. Arithmetic bias in the distribution of Mersenne
prime exponents and the divisor structure of $p-1$ [Электронный ресурс] // arXiv.
2026. arXiv:2603.08994. URL: \url{https://arxiv.org/abs/2603.08994}
(дата обращения: 27.08.2026).

\bibitem{jha2026} Jha A. The Poisson tail conjecture for primes in short intervals
[Электронный ресурс] // arXiv. 2026. arXiv:2605.23014.
URL: \url{https://arxiv.org/abs/2605.23014} (дата обращения: 27.08.2026).

\bibitem{gimps} PrimeNet [Электронный ресурс] / Great Internet Mersenne Prime
Search (GIMPS). Программное обеспечение Prime95/MPrime, версия~30.8.
URL: \url{https://www.mersenne.org/} (дата обращения: 27.08.2026;
отчёт \emph{Milestones} считан 27.08.2026 04:30~UTC).

\end{thebibliography}

% ── Романизированный список для международных баз цитирования ──────────────
% Все использованные источники англоязычны, поэтому транслитерация не
% требуется: список воспроизводит те же записи в международном формате.
\begin{otherlanguage}{english}
\section*{References}
\begin{enumerate}\itemsep2pt
\item Wagstaff S.~S., Jr. Divisors of Mersenne numbers.
\emph{Mathematics of Computation}, 1983, vol.~40, no.~161, pp.~385--397.
DOI: 10.1090/S0025-5718-1983-0679454-X

\item Pomerance C. On the distribution of pseudoprimes.
\emph{Mathematics of Computation}, 1981, vol.~37, no.~156, pp.~587--593.
DOI: 10.1090/S0025-5718-1981-0628717-0

\item Heuristics: deriving the Wagstaff Mersenne conjecture. \emph{The PrimePages:
prime number research and records}. Available at:
\url{https://t5k.org/mersenne/heuristic.html} (accessed 27.08.2026).

\item Dubner H. Generalized repunit primes.
\emph{Mathematics of Computation}, 1993, vol.~61, no.~204, pp.~927--930.
DOI: 10.1090/S0025-5718-1993-1185243-9

\item Bourdelais P. A generalized repunit conjecture. \emph{NMBRTHRY mailing
list}, 25 June 2009. Available at:
\url{https://listserv.nodak.edu/cgi-bin/wa.exe?A2=NMBRTHRY;417ab0d6.0906}
(accessed 25.08.2026).

\item OEIS Foundation Inc. \emph{The On-Line Encyclopedia of Integer
Sequences}. Available at: \url{https://oeis.org} (accessed 27.08.2026).

\item Wagstaff S.~S., Jr. Two conjectures about Mersenne primes.
\emph{Journal of Integer Sequences}, 2025, vol.~28, article 25.7.2.
Available at:
\url{https://cs.uwaterloo.ca/journals/JIS/VOL28/Wagstaff/wag3.html}
(accessed 25.08.2026).

\item Dominguez J. Arithmetic bias in the distribution of Mersenne prime exponents
and the divisor structure of $p-1$. \emph{arXiv}, 2026, arXiv:2603.08994. Available
at: \url{https://arxiv.org/abs/2603.08994} (accessed 27.08.2026).

\item Jha A. The Poisson tail conjecture for primes in short intervals.
\emph{arXiv}, 2026, arXiv:2605.23014. Available at:
\url{https://arxiv.org/abs/2605.23014} (accessed 27.08.2026).

\item Great Internet Mersenne Prime Search (GIMPS). \emph{PrimeNet}.
Prime95/MPrime software, version 30.8. Available at:
\url{https://www.mersenne.org/} (accessed 27.08.2026).
\end{enumerate}
\end{otherlanguage}

% ── Сведения об авторе ────────────────────────────────────────────────────
% п. 4.9 ГОСТ Р 7.0.7-2021: имя, отчество, фамилия полностью; организация,
% город, страна; электронный адрес БЕЗ слова «e-mail»; ORCID в форме URL.
% Элементы разделяют запятыми, точку в конце e-mail и ORCID не ставят.
\section*{Информация об авторе}
\noindent\textbf{Тимофей \textcolor{red}{[Отчество]} Докучаев} ---
\textcolor{red}{[учёная степень и звание при наличии; должность и наименование
организации с указанием подразделения --- если статья подаётся от организации]},
независимый исследователь, Иркутск, Россия,
\texttt{timofei93@yandex.ru}, \url{https://orcid.org/0009-0006-0510-5225}

\begin{otherlanguage}{english}
\section*{Information about the author}
\noindent\textbf{Timofei \textcolor{red}{[Patronymic]} Dokuchaev} ---
\textcolor{red}{[academic degree and title, if any; position and affiliation ---
if the article is submitted on behalf of an organization]},
independent researcher, Irkutsk, Russia,
\texttt{timofei93@yandex.ru}, \url{https://orcid.org/0009-0006-0510-5225}
\end{otherlanguage}

% п. 4.20.3: сведения об отсутствии конфликта интересов приводят в конце
% статьи ПОСЛЕ «Информации об авторе», на языке статьи и затем на английском.
\medskip
\noindent Автор заявляет об отсутствии конфликта интересов.

\begin{otherlanguage}{english}
\noindent The author declares no conflicts of interests.
\end{otherlanguage}

\end{document}
